\chapter{Заключение Тома VII: рангоизменяющий родитель}
\label{ch:v7-final-conclusion}

Том VII был открыт для проверки следующего сильного утверждения: может ли
одно физическое поле переменного ранга на общем аффинном, хиральном,
семейном и Real-носителе само вывести неустойчивость нулевого сектора,
выбрать физически допустимые направления и нелинейно завершить переход
устойчивым вакуумом. Только после этого разрешалось переходить к массам,
смешиваниям и именованию новых частиц.

Проверка завершена. Полного физического замыкания не получено, но впервые
построен один нормировочно устойчивый качественный класс действий,
действительно соединяющий ранговый запуск с положительным локальным
насыщением. Тем самым Том VII имеет смешанный итог: центральная
качественная задача решена на строгом амплитудно-тяжёлом срезе, а
количественная и феноменологическая части получили точные запреты.

\section{Исходные задачи и фактический результат}

\begin{enumerate}
 \item \textbf{Общий рангоизменяющий носитель --- пройдено.}
 После удаления повторного семейного множителя получен исправленный модуль
 $E_{\rm aff}\otimes\Lambda_{\rm ch}$ и связанные рёберный и цепной
 носители.

 \item \textbf{Стационарный неустойчивый нуль --- пройдено.}
 Hodge-момент даёт ровно семь отрицательных направлений; двадцать
 конкурирующих направлений имеют положительную щель.

 \item \textbf{Конечное нелинейное насыщение --- пройдено на срезе.}
 На 27-мерном вещественном срезе целевая точка стационарна и имеет
 сигнатуру $(0,0,27)$.

 \item \textbf{Один функционал запуска и насыщения --- пройдено
 качественно.} Рёберная Hodge-норма и полярная относительная кривизна
 собраны в $\mathcal S_\eta=\mathcal S_E+\eta\|\mathcal R_U\|^2$;
 представитель $\eta=1$ получается одним следом.

 \item \textbf{Единственная относительная метрика --- не пройдено.}
 Полная степень, Real-структура, Hodge-инволюция, кратности, полуслед,
 клиффордова степень, аффинный дефект и минимальная опора не фиксируют
 $\eta$.

 \item \textbf{Сохранение физической калибровочной группы --- частично.}
 Независимый Casimir- и графовый селекторы совпали и выделили бесцветный
 сектор; единый полный калибровочный след всего родителя не получен.

 \item \textbf{Семейные фазы и смешивания --- не пройдено.}
 Один цикл оставляет $U(3)/\operatorname{Ad}U(3)$; одна классовая функция
 не выбирает относительные семейные оси.

 \item \textbf{Масштаб и массовый спектр --- частично.}
 Одна масса калибрует линейный масштаб, но отношения полного спектра
 зависят от $\eta$, а квартика --- от $f_0$.

 \item \textbf{Полная пространственно-временная теория --- не пройдена.}
 Минимальный произведённый оператор даёт правильную форму кинетики и
 потенциала, но общий физический gauge--Hodge след не атрибутирован.

 \item \textbf{Частицы и наблюдаемые --- переход запрещён.}
 Массы, CKM/PMNS и новые частицы не объявляются, поскольку необходимые
 метрики и оси не выведены.
\end{enumerate}

\section{Что получилось}

\subsection{Ранг стал динамической переменной}

Исходный чисто кривизностный кандидат провалился, но хиральный
Hodge-функционал дал стационарный нуль и минимум меньшего ранга. После
исправления носителя физический ранг изменяется
\begin{equation}
 24\longrightarrow21,
\end{equation}
и оставляет три комплексные ядерные линии. Ранг и проектор здесь являются
следствиями поля, а не вставленными метками.

\subsection{Получен физически различающий запуск}

Цветосохраняющий квадратичный селектор и полный down--weak контроль дают
в начале
\begin{equation}
 (n_-,n_0,n_+)=(7,0,20).
\end{equation}
Это означает, что нулевой сектор не просто имеет отрицательную моду:
выбираются ровно семь требуемых корневых направлений, тогда как двадцать
ближайших конкурентов не запускаются.

\subsection{Найден устойчивый локальный конец перехода}

Полярная связывающая кривизна
\begin{equation}
 \mathcal R_U=C_tU-UC_s
\end{equation}
возникает как компонент квадрата трёхступенной нечётной цепи. Производная
по оператору числа удаляет диагональные возвраты и сохраняет именно эту
компоненту. В целевой точке полный амплитудно-тяжёлый гессиан равен
\begin{equation}
 (n_-,n_0,n_+)=(0,0,27).
\end{equation}

\subsection{Выделено нормировочно независимое ядро}

Главный положительный итог тома состоит не в особом числе $\eta=1$, а в
том, что вся область $\eta>0$ имеет одинаковую качественную картину:
\begin{equation}
 (7,0,20)_0\longrightarrow(0,0,27)_*.
 \label{eq:v7-final-qualitative-transition}
\end{equation}
Тем самым селектор и существование устойчивого локального насыщения не
являются следствием подгонки относительной массовой нормировки.

\section{Что не получилось}

\subsection{Не получен единственный количественный родитель}

Минимальная связывающая опора порождает фактор $M_{21}(\mathbb C)$, но
общая порождённая алгебра остаётся прямой суммой. В физической факторной
модели это препятствие имеет вид
\begin{equation}
 M_{22}(\mathbb C)\oplus M_{21}(\mathbb C),
\end{equation}
и оставляет один центральный относительный вес. Поэтому
\eqref{eq:v7-final-qualitative-transition} строго, а конкретные
собственные значения масс --- нет.

\subsection{Полный фазовый endpoint не изолирован}

Ранний 72-мерный гессиан обнаружил семейные кадровые моды. Циклическое и
лесное факторирования сократили их до одной голономии, но не вывели её
собственные фазы и оси. Высшие циклические характеры видят фазы, однако
нецентральный минимум требует свободного отношения моментов профиля, а
одна сопряжённо-инвариантная функция не выбирает CKM/PMNS.

\subsection{Масштаб и пространственно-временная нормировка не замкнуты}

Одна размерная калибровка фиксирует общий линейный масштаб и некоторые
условные отношения внутри чистого Hodge-сектора. Но полная массовая
матрица зависит от $\eta$, эффективная квартика имеет вид
$\lambda_E=\pi^2/f_0$, а общий физический след, одновременно фиксирующий
калибровочную кинетику и $f_0$, отсутствует.

\subsection{Физические частицы не выведены}

Том получил поля, проекторы, ядра, циклы и устойчивые моды. Но эти объекты
не получают автоматически названий известных или новых частиц. Для этого
нужны полный калибровочный модуль, квантовые числа, каноническая кинетика,
массовый оператор и наблюдаемые каналы. Эти требования не выполнены.

\section{Что было пересмотрено}

\begin{description}
 \item[Рангоизменяющий родитель.] Он существует не как единственная
 количественная теория, а как качественный универсальный класс действий.

 \item[Один общий след.] Он достаточен для предъявления одного допустимого
 представителя, но не доказывает единственность следа абстрактной прямой
 суммы.

 \item[Связывающая цепь.] Её физически активная опора имеет размерность
 $21$, а не $42$; средняя ступень является тождественным нулём кривизны.

 \item[Дефект $54-42=12$.] Совпадение с $\dim E_{\rm aff}$ оказалось
 артефактом нулевого дополнения, а не новым индексом.

 \item[Массовое отношение.] Отношение $\sqrt2$ сохраняется только внутри
 раннего чистого Hodge-сектора с общей кинетической метрикой. Оно не
 является спектром полного связанного родителя.

 \item[Смешивание.] Собственные фазы одного цикла не равны матрицам CKM или
 PMNS: для выбора относительных осей нужен второй некоммутирующий семейный
 тензор.
\end{description}

\section{Ответ на центральный вопрос}

\begin{theorem}[Ответ Тома VII]
В проверенной архитектуре существует общий исправленный носитель и
нормировочно устойчивый класс положительных действий, который на полном
исследованном амплитудно-тяжёлом срезе переводит стационарный нуль с семью
физическими отрицательными модами в строго устойчивую целевую точку. Тем
самым качественный рангоизменяющий переход получен. Однако единственная
массовая метрика, полный фазовый и пространственно-временной гессиан,
семейные смешивания и абсолютный масштаб не выведены. Поэтому полный
физический рангоизменяющий родитель и его спектр частиц не получены.
\end{theorem}

Это утверждение завершает поставленный различающий тест. Оно не является
универсальным запретом иной геометрии и не отменяет положительный
качественный механизм.

\section{Условия следующей программы}

Продолжение не должно состоять в новом выборе $\eta$ или функции
отсечения. Новая программа допустима только после предъявления одного из
трёх отсутствующих примитивов.

\subsection*{Новый примитив A: межсекторный коннектор}

Требуется физически типизированный оператор между рёберной и связывающей
опорами, делающий их проекторы нецентральными. Простое объявление полного
$M_{75}(\mathbb C)$ без такого оператора запрещено.

\subsection*{Новый примитив B: второй семейный тензор}

Для смешивания необходимо независимо вывести второй семейный оператор,
не коммутирующий с голономией первого цикла. Его параметры не могут быть
выбраны по CKM или PMNS.

\subsection*{Новый примитив C: общий калибровочно-пространственный след}

Один физический произведённый оператор должен одновременно фиксировать
кинетику, скалярный потенциал, калибровочные индексы, $f_0$ и
относительную Hodge-метрику. Только после этого допустимы абсолютные массы
и нелинейные амплитуды.

Каждый новый примитив обязан пройти Real-совместимость, калибровочную
типизацию, полный физический гессиан, ограниченность действия снизу и
слепую проверку без выбора параметров по целевым наблюдаемым.

\section{Итоговый статус}

\begin{protocol}
\begin{enumerate}
 \item исходная программа Тома VII: \textbf{проверена полностью};
 \item общий рангоизменяющий носитель: \textbf{получен после коррекции};
 \item стационарный неустойчивый нуль: \textbf{получен};
 \item селективный запуск $7+20$: \textbf{получен};
 \item устойчивое локальное насыщение на 27-мерном срезе:
 \textbf{получено};
 \item нормировочно независимый качественный класс: \textbf{получен};
 \item единственная массовая метрика: \textbf{не получена};
 \item фазовый и семейный endpoint: \textbf{не замкнут};
 \item общий калибровочно-пространственный след: \textbf{не получен};
 \item CKM/PMNS и абсолютные массы: \textbf{не выведены};
 \item новые физические частицы: \textbf{не заявлены};
 \item Том VII: \textbf{завершён как качественное замыкание и карта
 границы количественного родителя};
 \item следующая программа: \textbf{не открывается автоматически}.
\end{enumerate}
\end{protocol}